\subsection{\label{sec.gf.multivar}Multivariate FPSs}%

\[%
\begin{tabular}
[c]{||l||}\hline\hline
\textbf{TODO: This section needs more details.}\\\hline\hline
\end{tabular}
\]


Multivariate FPSs are just FPSs in several variables. Their theory is mostly
analogous to the theory of univariate FPSs (which is what we have been
studying above), but requires more subscripts. I will just discuss the main differences.

For instance, FPSs in two variables $x$ and $y$ have the form $\sum
_{i,j\in\mathbb{N}}a_{i,j}x^{i}y^{j}$, where the summation sign $\sum
_{i,j\in\mathbb{N}}$ means $\sum_{\left(  i,j\right)  \in\mathbb{N}^{2}}$ of
course. Formally, such an FPS is a family $\left(  a_{i,j}\right)  _{\left(
i,j\right)  \in\mathbb{N}^{2}}$ of elements of $K$. The indeterminates $x$ and
$y$ are defined by%
\[
x=\left(  \delta_{\left(  i,j\right)  ,\left(  1,0\right)  }\right)  _{\left(
i,j\right)  \in\mathbb{N}^{2}}\ \ \ \ \ \ \ \ \ \ \text{and}%
\ \ \ \ \ \ \ \ \ \ y=\left(  \delta_{\left(  i,j\right)  ,\left(  0,1\right)
}\right)  _{\left(  i,j\right)  \in\mathbb{N}^{2}}%
\]
(so that each indeterminate has exactly one coefficient equal to $1$, while
all other coefficients are $0$). The rules for addition, subtraction, scaling
and multiplication are essentially as they are for univariate FPSs, except
that now the indexing set is $\mathbb{N}^{2}$ instead of $\mathbb{N}$. For
example, multiplication of FPSs in $x$ and $y$ is defined by the formula%
\[
\left[  x^{n}y^{m}\right]  \left(  ab\right)  =\sum_{\substack{\left(
i,j\right)  ,\ \left(  k,\ell\right)  \in\mathbb{N}^{2};\\i+k=n;\\j+\ell
=m}}\left[  x^{i}y^{j}\right]  a\cdot\left[  x^{k}y^{\ell}\right]  b
\]
(for any two FPSs $a$ and $b$ and any $n,m\in\mathbb{N}$).

More generally, for any $k\in\mathbb{N}$, the FPSs in $k$ variables
$x_{1},x_{2},\ldots,x_{k}$ are defined to be the families $\left(
a_{\mathbf{i}}\right)  _{\mathbf{i}\in\mathbb{N}^{k}}$ of elements of $K$
indexed by $k$-tuples $\mathbf{i}=\left(  i_{1},i_{2},\ldots,i_{k}\right)
\in\mathbb{N}^{k}$. Addition, subtraction and scaling of such families is
defined entrywise. Multiplication is defined by the formula%
\[
\left[  \mathbf{x}^{\mathbf{n}}\right]  \left(  ab\right)  =\sum
_{\substack{\mathbf{i},\mathbf{j}\in\mathbb{N}^{k};\\\mathbf{i}+\mathbf{j}%
=\mathbf{n}}}\left[  \mathbf{x}^{\mathbf{i}}\right]  a\cdot\left[
\mathbf{x}^{\mathbf{j}}\right]  b
\]
(for any two FPSs $a$ and $b$ and any $\mathbf{n}\in\mathbb{N}^{k}$), where

\begin{itemize}
\item the sum $\mathbf{i}+\mathbf{j}$ means the entrywise sum of the
$k$-tuples $\mathbf{i}$ and $\mathbf{j}$ (that is, $\mathbf{i}+\mathbf{j}%
=\left(  i_{1}+j_{1},i_{2}+j_{2},\ldots,i_{k}+j_{k}\right)  $, where
$\mathbf{i}=\left(  i_{1},i_{2},\ldots,i_{k}\right)  $ and $\mathbf{j}=\left(
j_{1},j_{2},\ldots,j_{k}\right)  $);

\item if $\mathbf{m}\in\mathbb{N}^{k}$ and $h$ is an FPS in $x_{1}%
,x_{2},\ldots,x_{k}$, then $\left[  \mathbf{x}^{\mathbf{m}}\right]  h$ is the
$\mathbf{m}$-th entry of the family $h$. Just as in the univariate case, this
entry $\left[  \mathbf{x}^{\mathbf{m}}\right]  h$ is called the
\emph{coefficient of the monomial} $\mathbf{x}^{\mathbf{m}}:=x_{1}^{m_{1}%
}x_{2}^{m_{2}}\cdots x_{k}^{m_{k}}$ in $h$ (where $\mathbf{m}=\left(
m_{1},m_{2},\ldots,m_{k}\right)  $).
\end{itemize}

With this notation, the formula for $ab$ doesn't look any more complicated
than the analogous formula in the univariate case. The only difference is that
the monomials and the coefficients are now indexed not by the nonnegative
integers, but by the $k$-tuples of nonnegative integers instead.

The indeterminates $x_{1},x_{2},\ldots,x_{k}$ are defined by%
\[
x_{i}=\left(  \delta_{\mathbf{n},\left(  0,0,\ldots,0,1,0,0,\ldots,0\right)
}\right)  _{\mathbf{n}\in\mathbb{N}^{k}},
\]
where the tuple $\left(  0,0,\ldots,0,1,0,0,\ldots,0\right)  $ is a $k$-tuple
with a lone $1$ in its $i$-th position. Thus, if $\mathbf{m}=\left(
m_{1},m_{2},\ldots,m_{k}\right)  \in\mathbb{N}^{k}$, then%
\[
x_{1}^{m_{1}}x_{2}^{m_{2}}\cdots x_{k}^{m_{k}}=\left(  \delta_{\mathbf{n}%
,\mathbf{m}}\right)  _{\mathbf{n}\in\mathbb{N}^{k}},
\]
so that each FPS $f=\left(  f_{\mathbf{m}}\right)  _{\mathbf{m}\in
\mathbb{N}^{k}}$ satisfies%
\[
f=\sum_{\mathbf{m}=\left(  m_{1},m_{2},\ldots,m_{k}\right)  \in\mathbb{N}^{k}%
}f_{\mathbf{m}}x_{1}^{m_{1}}x_{2}^{m_{2}}\cdots x_{k}^{m_{k}}.
\]


Most of what we have said about FPSs in one variable applies similarly to FPSs
in multiple variables. The proofs are similar but more laborious due to the
need for subscripts. Instead of the derivative of an FPS, there are now $k$
derivatives (one for each variable); they are called \emph{partial
derivatives}. One needs to be somewhat careful with substitution -- e.g., one
cannot substitute non-commuting elements into a multivariate polynomial. For
example, you cannot substitute two non-commuting matrices $A$ and $B$ for $x$
and $y$ into the polynomial $xy$, at least not without sacrificing the rule
that a value of the product of two polynomials should be the product of their
values (since $\underbrace{x\left[  A,B\right]  }_{=A}\cdot
\underbrace{y\left[  A,B\right]  }_{=B}=AB$ would have to equal
$\underbrace{y\left[  A,B\right]  }_{=B}\cdot\underbrace{x\left[  A,B\right]
}_{=A}=BA$). But you can still substitute $k$ commuting elements for the $k$
indeterminates in a $k$-variable polynomial. You can also compose multivariate
FPSs as long as appropriate summability conditions are satisfied.

\begin{definition}
Let $k\in\mathbb{N}$. The $K$-algebra of all FPSs in $k$ variables
$x_{1},x_{2},\ldots,x_{k}$ over $K$ will be denoted by $K\left[  \left[
x_{1},x_{2},\ldots,x_{k}\right]  \right]  $.
\end{definition}

Sometimes we will use different names for our variables. For example, if we
work with $2$ variables, we will commonly call them $x$ and $y$ instead of
$x_{1}$ and $x_{2}$. Correspondingly, we will use the notation $K\left[
\left[  x,y\right]  \right]  $ (instead of $K\left[  \left[  x_{1}%
,x_{2}\right]  \right]  $) for the $K$-algebra of FPSs in these two variables.

Let me give an example of working with multivariate FPSs.

Let us work in $K\left[  \left[  x,y\right]  \right]  $. On the one hand, we
have%
\begin{align*}
\sum_{n,k\in\mathbb{N}}\dbinom{n}{k}x^{n}y^{k}  &  =\sum_{n\in\mathbb{N}%
}\ \ \sum_{k\in\mathbb{N}}\dbinom{n}{k}x^{n}y^{k}=\sum_{n\in\mathbb{N}}%
x^{n}\underbrace{\sum_{k\in\mathbb{N}}\dbinom{n}{k}y^{k}}_{\substack{=\left(
1+y\right)  ^{n}\\\text{(by the binomial formula)}}}\\
&  =\sum_{n\in\mathbb{N}}x^{n}\left(  1+y\right)  ^{n}=\sum_{n\in\mathbb{N}%
}\left(  x\left(  1+y\right)  \right)  ^{n}=\dfrac{1}{1-x\left(  1+y\right)
}\\
&  \ \ \ \ \ \ \ \ \ \ \left(
\begin{array}
[c]{c}%
\text{here, we have substituted }x\left(  1+y\right)  \text{ for }x\text{
in}\\
\text{the formula }\sum_{n\in\mathbb{N}}x^{n}=\dfrac{1}{1-x}\text{;}\\
\text{this is allowed since }x\left(  1+y\right)  \text{ has constant term }0
\end{array}
\right) \\
&  =\dfrac{1}{1-x}\cdot\dfrac{1}{1-\dfrac{x}{1-x}y}\ \ \ \ \ \ \ \ \ \ \left(
\text{easy to check by computation}\right) \\
&  =\dfrac{1}{1-x}\cdot\sum_{k\in\mathbb{N}}\left(  \dfrac{x}{1-x}y\right)
^{k}\\
&  \ \ \ \ \ \ \ \ \ \ \left(
\begin{array}
[c]{c}%
\text{here, we have substituted }\dfrac{x}{1-x}y\text{ for }x\\
\text{in the formula }\dfrac{1}{1-x}=\sum_{k\in\mathbb{N}}x^{k}%
\end{array}
\right) \\
&  =\dfrac{1}{1-x}\cdot\sum_{k\in\mathbb{N}}\dfrac{x^{k}}{\left(  1-x\right)
^{k}}y^{k}=\sum_{k\in\mathbb{N}}\dfrac{x^{k}}{\left(  1-x\right)  ^{k+1}}%
y^{k}.
\end{align*}
On the other hand, we have%
\[
\sum_{n,k\in\mathbb{N}}\dbinom{n}{k}x^{n}y^{k}=\sum_{k\in\mathbb{N}}%
\ \ \sum_{n\in\mathbb{N}}\dbinom{n}{k}x^{n}y^{k}=\sum_{k\in\mathbb{N}}\left(
\sum_{n\in\mathbb{N}}\dbinom{n}{k}x^{n}\right)  y^{k}.
\]
Comparing these two equalities, we find%
\begin{equation}
\sum_{k\in\mathbb{N}}\dfrac{x^{k}}{\left(  1-x\right)  ^{k+1}}y^{k}=\sum
_{k\in\mathbb{N}}\left(  \sum_{n\in\mathbb{N}}\dbinom{n}{k}x^{n}\right)
y^{k}. \label{eq.fps.mulvar.exa1.xyk}%
\end{equation}


Now, comparing coefficients in front of $x^{i}y^{k}$ in
(\ref{eq.fps.mulvar.exa1.xyk}), we conclude that%
\begin{equation}
\dfrac{x^{k}}{\left(  1-x\right)  ^{k+1}}=\sum_{n\in\mathbb{N}}\dbinom{n}%
{k}x^{n}\ \ \ \ \ \ \ \ \ \ \text{for each }k\in\mathbb{N}.
\label{eq.fps.mulvar.exa1.res1}%
\end{equation}
Let me explain in a bit more detail what I mean by \textquotedblleft comparing
coefficients in front of $x^{i}y^{k}$\textquotedblright. What we have used is
the following simple fact:

\begin{proposition}
\label{prop.fps.mulvar.comp-y-coeff}Let $f_{0},f_{1},f_{2},\ldots$ and
$g_{0},g_{1},g_{2},\ldots$ be FPSs in a single variable $x$ such that%
\begin{equation}
\sum_{k\in\mathbb{N}}f_{k}y^{k}=\sum_{k\in\mathbb{N}}g_{k}y^{k}%
\ \ \ \ \ \ \ \ \ \ \text{in }K\left[  \left[  x,y\right]  \right]  .
\label{eq.prop.fps.mulvar.comp-y-coeff.ass}%
\end{equation}
Then, $f_{k}=g_{k}$ for each $k\in\mathbb{N}$.
\end{proposition}

\begin{proof}
For each $k\in\mathbb{N}$, let us write the two FPSs $f_{k}$ and $g_{k}$ as
$f_{k}=\sum_{n\in\mathbb{N}}f_{k,n}x^{n}$ and $g_{k}=\sum_{n\in\mathbb{N}%
}g_{k,n}x^{n}$ with $f_{k,n},g_{k,n}\in K$. Then, the equality
(\ref{eq.prop.fps.mulvar.comp-y-coeff.ass}) can be rewritten as
\[
\sum_{k\in\mathbb{N}}\ \ \sum_{n\in\mathbb{N}}f_{k,n}x^{n}y^{k}=\sum
_{k\in\mathbb{N}}\ \ \sum_{n\in\mathbb{N}}g_{k,n}x^{n}y^{k}.
\]
Now, comparing coefficients in front of $x^{n}y^{k}$ in this equality, we
obtain $f_{k,n}=g_{k,n}$ for each $k,n\in\mathbb{N}$. Therefore, $f_{k}=g_{k}$
for each $k\in\mathbb{N}$. This proves Proposition
\ref{prop.fps.mulvar.comp-y-coeff}.
\end{proof}

Now, because of (\ref{eq.fps.mulvar.exa1.xyk}), we can apply Proposition
\ref{prop.fps.mulvar.comp-y-coeff} to $f_{k}=\dfrac{x^{k}}{\left(  1-x\right)
^{k+1}}$ and $g_{k}=\sum_{n\in\mathbb{N}}\dbinom{n}{k}x^{n}$. Thus, we obtain
the equality%
\[
\dfrac{x^{k}}{\left(  1-x\right)  ^{k+1}}=\sum_{n\in\mathbb{N}}\dbinom{n}%
{k}x^{n}\ \ \ \ \ \ \ \ \ \ \text{for each }k\in\mathbb{N}.
\]
This is an equality between \textbf{univariate} FPSs, even though we have
obtained it by manipulating \textbf{bivariate} FPSs (i.e., FPSs in two
variables $x$ and $y$). As a homework exercise (Exercise
\ref{exe.fps.sum-n-choose-k-xn-elementary}), you can prove this equality in a
more elementary, univariate way. But the idea to introduce extra variables (in
our case, in order to discover an equality) is highly useful in many
situations (some of which we will see later in this course).

\section{\label{chap.pars}Integer partitions and $q$-binomial coefficients}

We have previously counted compositions of an $n\in\mathbb{N}$. These are
(roughly speaking) ways to write $n$ as a sum of finitely many positive
integers, where the order matters. For example, $3=3=2+1=1+2=1+1+1$, so $3$
has $4$ compositions. Formally, compositions are tuples of positive integers.

Now, let us disregard the order. There are two ways to make this rigorous:
either we replace tuples by multisets, or we require the tuples to be weakly
decreasing. These result in the same count, but we will use the 2nd way, just
because tuples are easier to work with than multisets. This will lead to the
notion of \emph{integer partitions}.
